\documentclass[10pt]{article}
\usepackage[margin=20mm]{geometry}
\usepackage{amsmath,amssymb,mathtools,fontspec,unicode-math,xurl,hyperref}
\IfFontExistsTF{DejaVu Serif}{\setmainfont{DejaVu Serif}}{\setmainfont{Latin Modern Roman}}
\setmathfont{latinmodern-math.otf}
\setlength{\parindent}{0pt}\setlength{\parskip}{6pt}\sloppy
\title{Initial review of regularity of the singular set in optimal transportation}
\author{}
\date{26 October 2026}
\begin{document}
\maketitle
Abstract: The checked scope comprises localization underlying the first regularity theorem, early obliqueness estimates, and the density argument leading to higher regularity. A literal affine-invariance equality is false, but the precise bound required follows directly from the existing chord endpoints. The invalid inference occurs in the core localization proof and is recorded as a Major proof-step error. The supplied repair establishes the required estimate; no counterexample to the regularity theorem is asserted. The long obliqueness analysis and external regularity theorems remain unaudited dependencies. Scientific acceptance is pending.\par
\section*{Identity and exact scope}
The reviewed source is Shibing Chen and Jiakun Liu, Regularity of singular set in optimal transportation, arXiv:2210.13841v4, revised 20 July 2025 [1]. It has 43 PDF pages; PDF and printed page numbers agree. The bound file is `papers/0092/0092.pdf`. The associated journal record is Inventiones mathematicae 242 (2025), 1–44, DOI `10.1007/s00222-025-01353-w`, published 30 July 2025 [2]. The journal version of record was not acquired; no equivalence with that version is asserted.\par
Pages 1–19 and 39–42 were read in extracted text and actually viewed as page images (23 pages). The focus was Theorems 1.1–1.2, Lemmas 3.1–3.5, the beginning of Section 4 through Lemma 4.3, and Lemma 5.2. The statement and beginning of Lemma 4.4 were read, but its complete proof was not reviewed. Pages 20–38, containing much of the decisive obliqueness analysis, and page 43 were not scientifically inspected. Indexing the entire PDF is not counted as reviewing it.\par
\section*{Hypotheses, quantifiers, and dependencies}
The source and target are bounded open sets with probability densities bounded above and below by positive constants. The target has two convex components separated by a hyperplane. The source is not assumed convex. The convex transport potential \(u\) has branches \(u_i\), defined by restricted Legendre duality, and the interface is \[{\mathcal{F}}=\{u_1=u_2\}\cap\Omega .\] The claims are local inside \(\Omega\), including closures of branch regions at \(\mathcal F\), not assertions of regularity at an arbitrary exterior source boundary. Theorem 1.1 gives \(C^{1,\beta}\) regularity for some \(\beta\in(0,1)\). Theorem 1.2 additionally assumes \(C^2\) uniformly convex target components, \(f\in C^\alpha_{\mathrm{loc}}(\Omega)\) and \(g\in C^\alpha(\overline{\Omega^*})\), with \(0<\alpha<1\), and concludes \(C^{2,\alpha}\) regularity.\par
Centered and tangent sections are different objects. Their comparison in (4.7) is needed in the width estimates. Lemma 3.5 gives Alexandrov measure inequalities, not merely almost-everywhere determinant bounds.\par
The checked chain is: compactness of the zero set and exposed points give localized sections (Lemmas 3.1–3.3); differentiability puts nearby subgradients inside the source (Lemma 3.4); images of exterior target points then lie in the Lipschitz interface, a null set, proving absence of exterior Alexandrov mass (Lemma 3.5). Quantitative convexity and Legendre duality give branch \(C^{1,\beta}\) estimates. The nonzero gradient jump in (1.5) determines the interface normal. Existing results cited as [4], [5], [7], and [14] in the paper remain external dependencies.\par
\section*{Problem 1: an invalid affine-invariance step in Lemma 3.4}
\subsection*{Exact location and geometric check}
Equation reference: bad. On p. 13, the unnumbered display immediately before (3.20) asserts \[
 \frac{\operatorname{dist}(Lp,\partial LS)}{\operatorname{diam}(LS)}
 =
 \frac{\operatorname{dist}(p,\partial S)}{\operatorname{diam}(S)}
 \leq\frac{t_2}{t_1}\leq\frac1M ,\] where \(S=S^c_{h'}[v](y_{t_0})\) and \(L\) is a John normalization. The same equality occurs in the bound author TeX, lines 597–598. Affine maps preserve ratios along a fixed line, not this Euclidean boundary-distance/diameter quotient.\par
Let \[S=\{(x_1,x_2):x_1^2/4+x_2^2<1\},\quad
 L=\operatorname{diag}(1/2,1),\quad p=(0,1-\delta),\quad 0<\delta<1.\] This bounded convex centered section shape maps to the unit disk, with \(L0=0\). The distance from \(p\) to the ellipse boundary is \(\delta\): parametrizing a boundary point as \((2\cos\theta,\sin\theta)\) gives squared distance \[4-3s^2-2(1-\delta)s+(1-\delta)^2,\qquad s=\sin\theta\in[-1,1].\] This concave polynomial has its minimum at an endpoint; \(s=1\) gives \(\delta^2\), smaller than \((2-\delta)^2\). The original quotient is therefore \(\delta/4\) and the transformed quotient is \(\delta/2\). At \(\delta=1/100\), these are exactly \(1/400\) and \(1/200\).\par
This example tests only the affine geometric identity. It does not purport to satisfy every transport assumption of Theorems 1.1–1.2 and is not a counterexample to either theorem.\par
\subsection*{Complete replacement and impact}
Equation reference: good. The required (3.20) is nevertheless valid using the endpoints already constructed on pp. 11–12. Write \(Lx=Ax+b\) and \[q_+=p+t_2e,\quad q_-=p-t_1e\in\partial S,\qquad t_1,t_2>0 .\] Then \(\operatorname{dist}(Lp,\partial LS)\leq t_2|Ae|\) and \(\operatorname{diam}(LS)\geq(t_1+t_2)|Ae|\). Consequently \[
 \frac{\operatorname{dist}(Lp,\partial LS)}{\operatorname{diam}(LS)}
 \leq\frac{t_2}{t_1+t_2}
 \leq\frac{t_2}{t_1}\leq\frac1M .\] Since (3.19) gives \(\operatorname{diam}(LS)\leq2C_n\), this yields (3.20) after changing the dimensional constant. No condition-number bound or added transport hypothesis is needed. The ellipse also has the relevant chord parameters \(t_0=1-\delta\), \(t_1=2-\delta\), \(t_2=\delta\).\par
The comparison of (3.22) and (3.23) requires only this \(C_n/M\) bound; its large-\(M\) contradiction survives. Replace the literal equality by [good] (\#good). The paper explicitly invokes ordinary affine invariance of the Euclidean distance ratio; John normalization in (3.19) imposes no similarity condition and no pullback metric is defined. Chord-ratio invariance proves the replacement inequality, not the displayed equality.\par
This is a Major error in a core proof step. The false equality is used to obtain (3.20), hence (3.23), the contradiction establishing Lemma 3.4, and then (3.24), Lemma 3.5 and (3.25) in the proof of Theorem 1.1. The complete repair above uses the existing endpoints and adds no transport hypothesis. The geometric counterexample disproves the invoked affine principle; it is not a constructed transport solution satisfying every hypothesis of the regularity theorem. Neither theorem falsity nor author adoption is asserted.\par
\section*{Further checks and counterevidence}
\subsection*{Localization}
In Lemma 3.1, a ray in the zero set forces subgradients into a half-space by monotonicity; the positive source density near the normalized interior point rules this out. Lemmas 3.2–3.3 use exposedness and compact section limits. In Lemma 3.5, (3.24) puts an exterior target point’s subgradient in \(\Omega\). An interior branch point would send that gradient back into the corresponding target component by strict convexity and duality. The remaining subgradients lie in \(\mathcal F\), whose Lipschitz regularity makes its Lebesgue measure zero. The local ball avoids the other target component. No source convexity is assumed.\par
\subsection*{Small exponents}
For Lemma 4.2, the volume comparison on p. 18 gives \[d^{3/2}h^{(n-2)(1/2+\varepsilon)}
 \lesssim h^{n/2},\qquad
 d\lesssim h^{2/3-2(n-2)\varepsilon/3}.\] Renaming an arbitrarily small positive \(\varepsilon\) gives the stated exponent. In Lemma 4.3, the Legendre conjugate of a growth power \(3+\eta\) has power \[\frac{3+\eta}{2+\eta}=\frac32-\frac{\eta}{2(2+\eta)}.\] The stated \(3/2-\varepsilon\) loss is legitimate. The nonzero normal jump near the interface is needed to pass from gradient bounds to the graph estimate.\par
In Lemma 5.2, choose \(\varepsilon=\beta/[8(1+\beta)]\). Then \[(1/2-\varepsilon)(1+\beta)-(1/2+\beta/4)=\beta/8>0,\qquad
 \beta/4-\varepsilon>0.\] The graph-height threshold is asymptotically smaller than the section’s positive axial radius. The convex-segment argument places the upper cap inside \(\Omega_1\), and John normalization gives a fixed positive volume fraction for small \(h\). Convexity of \(\Omega_1\) is not needed.\par
\subsection*{Forward references and higher regularity}
The use of Lemma 5.4 in Case III on p. 39 is not by itself evidence of circularity: the relevant branch is already oblique in that case, and the estimates of Section 5 are branchwise once that branch’s obliqueness and the earlier interface regularity are available. A complete audit of Case I/II dependencies was not performed.\par
For care with the displayed powers there, conjugating an upper bound of exponent \(2-\eta\) yields a lower growth exponent \((2-\eta)/(1-\eta)>2\), rather than the same \(2-\eta\). The weaker resulting lower bound still dominates a power \(3+\varepsilon\) for sufficiently small \(\eta\), which is the comparison needed in that paragraph. This is not a transport counterexample and is not promoted to a confirmed substantive defect. The \(C^{2,\alpha}\) perturbation step on p. 42 invokes [10, Section 4]; that external proof and its full transfer were not independently audited.\par
\section*{Outcome and remaining work}
One Major proof-step error is established within the inspected scope: the affine-invariance assertion in Lemma 3.4. Its required replacement bound is proved, preserving the localization argument after correction. A targeted primary-source check found arXiv v4 as the listed latest version and the journal publication record; no article-specific correction was located in that limited search. This is not proof that no correction exists.\par
Question | Recorded disposition\par
Issue judgment | One Major error in the core localization proof; explicit repair verified.\par
Replacement | Analytically verifies (3.20), with an inessential dimensional constant change.\par
Adoption | Not requested or established.\par
Remaining work | Full obliqueness proof, external prerequisites and perturbation transfer, journal-version comparison.\par
Acceptance | Pending; no human review or scientific acceptance recorded.\par
\section*{Supporting checks}
It confirms the rational ellipse values and symbolic exponent margins; these support, but do not replace, the analytic arguments. Saved outputs are `evidence/independent\_checks.json` and `evidence/static\_source\_check.json`. Author source was read statically only; no author code was run. The available TeX and figures are typesetting assets, not a missing computational experiment.\par
> 00005TOKEN\par
The inspected author TeX SHA-256 is\par
> 00006TOKEN\par
The broad review, calculations and static inspection are retained initial-review work, identified in the preserved provenance. During focused independent acceptance, the current executor and a bounded read-only native helper rechecked the existing affine-invariance issue; the executor corrected this report, compiled its exact TeX and checked the changed pages. The main dispatcher supplies scheduling and attributed UI observations only. Human reviews remain empty and the scientific decision remains pending.\par
2 S. Chen and J. Liu, Regularity of singular set in optimal transportation, arXiv:2210.13841v4, 20 July 2025. https://arxiv.org/abs/2210.13841v4. Bound PDF and author TeX consulted. S. Chen and J. Liu, same title, Invent. Math. 242 (2025), 1–44, DOI `10.1007/s00222-025-01353-w`. Publisher metadata consulted 26 October 2026; journal full text not reviewed.\par
\end{document}
